\chapter{The grinding circuit and the comparison}
\label{ch:mill}
% Purpose (THESIS_GUIDE.md, Ch. 5): everything the reader needs to judge the
% mill comparison as fair: plant, loops, interface, scenarios, scores, the
% TD-MPC and how it is run, the noise models. State the choices and who agreed
% them; the fairness discussion goes to the Discussion. Target 8-10 pages.
\TODO{Introduction (one paragraph): what this chapter fixes, and that the
choices were agreed with Weir (Kaige) before the protocol runs.}

\section{Simulator and regulatory layer}
\label{sec:mill-plant}
\paragraph{Simulator.} The industrial case uses Weir's dynamic simulator of a
secondary ball-mill circuit in reverse configuration, written in JAX. Fresh
slurry enters the sump, a pump feeds the sump slurry to a cluster of
hydrocyclones, the cyclone overflow is the product, whose particle size P80 is
the quality variable, and the underflow returns through the ball mill to the
sump. Water is added at the sump and at the mill inlet. The circulating load
is the ratio of underflow solids to fresh-feed solids. The experiment uses the
setting of Weir's TD-MPC environment: explicit Euler integration with a 10~s
scan. Episodes last two hours and start from a warm start, the circuit
settled for three hours at the nominal recipe.
\TODO{Flowsheet figure (new, drawio): mill, sump, pump, cyclones, the five
setpoints and the loops.}

\paragraph{Regulatory layer.} The lowest layer is Weir's own controller chain
for the circuit, used unchanged: a feed generator, through which the
disturbances enter, and five PID loops, all sampled every second. The cyclone
pressure is controlled by the pump flow setpoint, with an override on the
sump level near the edges of its range; the pump flow by the pump speed; the
sump density by the setpoint of the sump water flow, which has its own inner
flow loop; and the mill water flow by its valve. The number of operating
cyclones and the mill speed are applied directly.
\TODO{Mention that the cyclone-valve loop stays off (redundant per Weir).}

\section{Operating point}
\label{sec:mill-operating-point}
The nominal feed is the centre of
the TD-MPC's training range: 1\,250~m$^3$/h of slurry, ore density
2.60~t/m$^3$ and 41.5\% solids by volume. The nominal recipe is 9 cyclones,
11~psi, a density setpoint of 1.697~t/m$^3$, 50~t/h of mill water and 13~rpm,
which gives a sump level of 75\% and a P80 of 188~$\mu$m.
\TODO{One or two sentences: the simulator's shipped recipe is not an
operating point (the sump drifts out of the band and the level override
limit-cycles; mill \texttt{docs/phase0\_findings.md},
\texttt{docs/remeasure\_mpc\_interface.md}). Phase~0 found this, and it
affects every controller. Only one operating point is studied
(delimitation).}

\section{Supervisor interface}
\label{sec:mill-interface}
The supervisor sets the same five
setpoints as the TD-MPC, within the same limits: the number of operating
cyclones (7--9), the cyclone pressure setpoint (9--13~psi), the sump density
setpoint (1.60--1.80~t/m$^3$), the mill water addition (30--120~t/h) and the
mill speed setpoint (10--15~rpm). It sees the TD-MPC's 30 sensors
(fresh feed, mill, sump, pump and cyclone measurements and the circulating
load) and may ask for up to eight more simulator signals. The objectives are
the P80 range 180--200~$\mu$m, the sump
level band 70--80\%, the safety bounds of the sump level, 20 and 95\%, at
which the simulator stops, and the circulating-load band 400--600\%.
\TODO{Decision interval: 300~s for the main comparison (Weir, 2026-11-03: a
recommended operator cadence), 60 and 900~s as variants. Window length and
sampling at 300~s. Wrapper move limits per decision (one cyclone, 1~psi,
0.05~t/m$^3$, 20~t/h, 1~rpm; check \texttt{configs/supervisor.yaml}, no
longer provisional). Measurements are noise-free in the main comparison.}
% Table: setpoint, limits, wrapper step per decision.

\section{Scenarios and batteries}
\label{sec:mill-scenarios}
Disturbances act on
the feed flow ($\pm$2--8\% in the development range), the ore density
($\pm$0.03--0.10~t/m$^3$) and the solids content ($\pm$0.5--1.5 percentage
points), so that the development range is the TD-MPC's training range. They
start after 5--20 minutes, as steps, ramps over 10--30 minutes or sines with
periods of 20--60 minutes. The nominal scenarios step the P80 target by
2.5--5\% and back.
\TODO{Batteries for training: 2 scenarios per cell, 28 development, 28
validation and 24 beyond-range (seeds 1, 2, 3). Beyond range: 1.25--2 times
the development half-widths, with explicit values (Kaige asked;
\texttt{llm\_supervisor/scenarios.py}), e.g.\ feed 1050--1150 or
1375--1450~m$^3$/h. The sealed test battery (216 scenarios, five parts;
Section~\ref{sec:protocol}). One table with all ranges per battery.}
% Table: disturbance kind x {development, beyond, sealed amplitude levels}.

\section{Score}
\label{sec:mill-score}
\TODO{The multi-objective score as an equation with one term per objective, so that the results can be read from several perspectives
(P80, level, circulating load, safety, energy). Weights frozen by the protocol.}
% Weights: score_weights in configs/supervisor.yaml (mill repo).
% J_s = 5 * (P80 outside 180-200 um, um.h)
%     + 0.1 * (s with level outside 70-80 %) + 0.1 * (s with CL outside 400-600 %)
%     + 1 * (s outside the safety bounds, incl. the rest of a stopped episode)
%     + 10 * (energy per tonne of ore, kWh/t)
%     + 0.02 * (|CL - 500 %|, %.h)
%     + 1000 * (exceptions)
% P80 IAE and setpoint travel are reported, not scored.
\TODO{The two decisions behind it: setpoint travel removed from the score
(mill \texttt{docs/travel\_weight\_analysis.md}: it was 80--90\% of the score
and drove deadband tuning); energy added (\texttt{docs/energy\_sweep.md}:
mill speed is the energy lever, 12~rpm saves about 6\% kWh/t at nominal
feed). Then the two derived measures, both descriptive only: the
\emph{control penalty} (score minus its energy term; energy is about 90\% of a
good controller's score and moves mostly with the feed) and the
\emph{normalised score} (1 = every sample in spec at the lowest sustainable
mill power, 0 after a terminal state; \texttt{docs/normalised\_score.md};
decided 2026-11-06: reported beside the score, training unchanged).}

\section{The TD-MPC}
\label{sec:mill-tdmpc}
\TODO{Introduce the model: \texttt{autumn-night-41}, the model deployed in
Weir's backend, trained by Weir on the same simulator setting with the feed
held constant over each episode within the development range; latent model,
planning the five setpoints over a 40-minute horizon; same
\texttt{tasks.yml}; its own planner configuration. Training cost and data if
Weir can give them (for the cost comparison in the Discussion).}

Weir's trained model is run in the same
harness as the LLM supervisors, on the same scenarios, through the same safety
wrapper and with the same score. At each decision it receives the last
15~minutes of the 30 sensors sampled once a minute, together with the P80
target, and returns the five setpoints. It was trained to hold each action for
300~s, so it is evaluated at its own interval of 300~s and, as variants, at
60~s, where it replans every minute and applies the first step of each plan,
and at 900~s.
Its planner is reset before every episode. Two differences between the
comparators are deliberate and reported rather than removed. First, the
TD-MPC was trained with the feed held constant over each episode, while most
scenarios vary the feed as ramps and sines, so the results are also split
into steady and step scenarios against ramp and sine scenarios. Second, the
TD-MPC maximises its own reward, not the score of this thesis: a product of
bell functions around the P80 target (weight 0.5), a sump level of 75\% (1.0)
and a circulating load of 500\% (0.25), plus a term in the net mill power and
a penalty near the sump limits. Every episode is therefore scored twice, with
the score and with the TD-MPC's reward per 10~s sample averaged over the
episode, where samples after the simulator stops count as zero.

\section{Baselines}
\label{sec:mill-baselines}
Besides the fixed recipe, a decentralised PI supervisor serves
as a simple baseline: it trims the density setpoint from the P80 error and the
pressure setpoint from the sump level.
\TODO{Its gains and how they were chosen; run at the same decision interval
as the others.}

\section{Noise models}
\label{sec:mill-noise}
% RQ2 must be quantified along noise categories (2026-11-19 meeting): organise
% the models by what the noise acts on (measurement vs process/feed), its
% amplitude, and its temporal structure. Parameters in Appendix app:noise.
\TODO{Introduce noise as its own axis of RQ2, organised in categories:}
% - Category (where it enters): measurement noise on what the controllers see
%   (scores stay on the true values) vs. process noise on the feed.
% - Amplitude: e.g. P80 sigma 5 um; levels to sweep if amplitude is a factor.
% - Temporal structure: white; coloured AR(1) (tau 5 min) plus drift (tau 1 h);
%   sample-and-hold (analyser every 120 s); feed fluctuations over the episode.
% The models:
% - P80 white, sigma 5 um (Kaige's variant in tasks.yml).
% - p80_analyser: coloured + drift, sigma 5 um total, read every 120 s and held.
% - plant: analyser P80 plus guessed instrument accuracies on every signal
%   (flows ~1 %, densities 0.01, level 0.5 %, pressure 0.15 psi).
% - feed15: real fresh-feed fluctuations, 15 % of flow peak-to-peak over 2 h.
% - Combinations (feed + measurement).
% Say which sizes are guesses (measurement sizes; feed sizes for ore density
% and solids await Weir). Draws are paired across controllers; noise-trained
% runs are tested on fresh draws (--noise-seed-offset).
% A table: model x category x amplitude x temporal structure x signals.
